\documentclass[../main.tex]{subfiles}
\begin{document}

\section{Problema Cal-cule (Lot 2025)}
\label{problem:cal-cule}

\href{https://kilonova.ro/problems/3820}{enunț}
$\bullet$
\hyperref[code:cal-cule]{surse}

\subsection{Enunțul}

Să clarificăm mai întîi enunțul problemei (uneori este necesar).

\begin{itemize}
  \item Există un șir secret de $n \leq \np{1000}$ elemente.

  \item La intrare primim $n$ și un alt număr $k$.

  \item Trebuie să aflăm histograma (distribuția frecvențelor) din șirul secret.

  \item Putem face interogări despre histograma oricărui subșir, specificînd indicii elementelor din subșirul dorit.

  \item \textbf{Dar} nu primim ca răspuns histograma, ceea ce ar fi prea simplu. În schimb, comisia trece histograma printr-o funcție pe care tot noi o implementăm, care reduce o histogramă la o valoare între $0$ și $k-1$.
\end{itemize}

Concret,

\begin{itemize}
  \item Trebuie să implementăm funcția \ccode{descuie_grajdul}.

  \item Din \ccode{descuie_grajdul} putem apela funcția comisiei, \ccode{intreaba_cal}, trimițîndu-i indicii din subșirul dorit.

  \item \ccode{intreaba_cal} calculează histograma subșirului, iar cu acea histogramă apelează funcția noastră, \ccode{dreseaza_cal}.

  \item Funcția \ccode{dreseaza_cal} trebuie să reducă orice vector la o valoare între $0$ și $k-1$.
\end{itemize}

Restricțiile includ un tabel de valori posibile pentru $k$. Vedem că, cu cît $k$ este mai mare, cu atît limita pentru numărul de interogări scade (firește).

\subsection{O soluție în \texorpdfstring{$\bigoh(n \log n)$}{O(n log n)}}

Prima mea idee a fost să aflu efectiv vectorul. Desigur, nu putem afla exact valorile, dar putem afla clasele de echivalență (pe ce poziții apar valori egale). De exemplu, este relevant că $v[0] = v[2] = v[7]$, dar nu este relevantă valoarea efectivă.

Să considerăm că am aflat clasele pe pozițiile $0, 1, \dots, i-1$ și dorim să aflăm clasa poziției $i$, adică dorim să aflăm cu ce poziție anterioară este egal $v[i]$ sau dacă $v[i]$ este o valoare încă nemaiîntîlnită. Putem căuta binar această informație astfel. Să colectăm indicii primei apariții a fiecărei valori distincte. De exemplu, dacă $i = 9$ și dacă știm că

\begin{itemize}
  \item $v[0] = v[2] = v[7]$
  \item $v[1] = v[3] = v[4] = v[6]$
  \item $v[5] = v[8]$
\end{itemize}

atunci vom colecta indicii 0, 1 și 5. Subliniem că elementele $v[0]$, $v[1]$ și $v[5]$ sînt distincte. Cu alte cuvinte, vectorul de frecvențe pentru subșirul $(v[0], v[1], v[5])$ va fi (1, 1, 1). Acum, dacă adăugăm elementul $v[9]$ la acest subșir apar două posibilități:

\begin{enumerate}
  \item Fie vectorul de frecvențe rezultat este (1, 1, 1, 1), caz în care $v[9]$ este o valoare nouă, diferită de $v[0 \dots 8]$.

  \item Fie vectorul de frecvențe rezultat este (1, 1, 2), caz în care $v[9]$ este egal cu unul dintre elementele anterioare.
\end{enumerate}

Astfel, putem căuta binar cel mai scurt prefix al subșirului format din indicii primelor apariții care, împreună cu poziția $v[i]$, produce o frecvență de 2. Pentru aceasta, vom face \ccode{dreseaza_cal} să returneze 0 sau 1 după cum vectorul de frecvențe se termină sau nu cu o frecvență 2.

Această soluție ajunge la \np{8000} de interogări, iar principalul său neajuns este că \ccode{dreseaza_cal} returnează doar două valori posibile, indiferent de valoarea lui $k$. Cu alte cuvinte, soluția nu încearcă să partiționeze spațiul de posibilități suficient de rapid.

\subsection{Soluția optimă}

Eu nu văd niciun mod de a împinge soluția anterioară mai departe. Prin natura ei, căutarea binară utilizează doar două valori din cele $k$, pentru că pune întrebări de tip da/nu.

În loc de asta, ne mulțumim să baleiem vectorul și să aflăm la fiecare pas histograma (distribuția frecvențelor). De exemplu, dacă primele 5 elemente sînt distincte, vom avea 5 valori cu frecvența 1. Dacă al 6-lea element este egal cu unul anterior, nu insistăm să aflăm \textbf{cu care}. Ne interesează doar noua histogramă: 4 frecvențe de 1 și o frecvență de 2.

Un indiciu valoros este că, pentru $k = \np{1000}$, avem doar 999 de interogări. Cu alte cuvinte, pentru fiecare poziție de la 1 la $n-1$ trebuie să facem exact o interogare. Atunci, haideți să presupunem că am pus întrebări pentru elementele $0, 1, \dots, i-1$ și să gîndim soluția prin prisma elementului $i$. Cînd îl procesăm, exact o frecvență va crește cu exact o unitate față de prefixul anterior. Dacă elementul este la prima apariție, un 0 devine 1. Altfel un 1 devine 2 sau un 2 devine 3 etc. Cum decidem? Am vrea să facem o interogare astfel încît toate cele $k$ răspunsuri posibile să fie relevante. De exemplu, răspunsul 0 să indice că o frecvență 0 devine 1, răspunsul 1 să indice că o frecvență 1 devine 2 etc.

Ideal, am putea dresa calul să raporteze numărul de perechi egale în prefixul interogat. De exemplu, șirul $(7, 8, 8, 9, 9, 9)$ conține 4 perechi (perechea 8-8 și 3 perechi 9-9). Dacă extindem șirul la $(7, 8, 8, 9, 9, 9, 8)$, el va conține 6 perechi, cu 2 mai multe decît anteriorul, ceea ce arată că o frecvență a crescut de la 2 la 3. Așadar, de acum încolo cînd spunem „răspuns” ne referim la perechile formate de elementul curent, calculate ca diferența dintre răspunsul dat de cal și numărul de perechi cunoscute din interogarea pe care noi am format-o.

Partea bună este că nu avem nevoie să știm șirul ca să aflăm numărul de perechi, ci doar frecvențele. Deci \ccode{dreseaza_cal} doar însumează $f \cdot (f-1)/2$ pentru fiecare frecvență $f$. Dar trebuie clarificate multe aspecte.

\subsubsection*{Cazul cu puține valori distincte}

Întîi, suma frecvențelor va depăși $k$. Cînd $k=n$, putem opera modulo $k$, deoarece la orice pas creșterea numărului de perechi nu va depăși $n-1$. Dar, cînd $k$ este mic, nu putem diferenția răspunsul. De exemplu, dacă $k=4$ și creșterea este 1, atunci s-ar putea ca o frecvență de 1 să fi devenit 2 sau una de 5 să fi devenit 6 etc. Putem reduce confuzia dacă numărul de elemente distincte este cel mult $k-1$. De exemplu, să considerăm următorul șir cu $k=4$, frecvențele 1, 5, 9 și elementul deocamdată necunoscut $x$:

\begin{verbatim}
8 13 13 13 13 13 20 20 20 20 20 20 20 20 20 x
\end{verbatim}

Vom construi interogarea astfel:

\begin{itemize}
  \item Elementul cu frecvența 1 îl luăm o dată.
  \item Elementul cu frecvența 5 îl luăm de două ori.
  \item Elementul cu frecvența 9 îl luăm de trei ori.
  \item Îl adăugăm și pe $x$.
\end{itemize}

Rezultă interogarea 8 13 13 20 20 20 $x$ (reamintesc că vom trimite poziții, nu valori, căci acesta este protocolul). Interpretăm răspunsul astfel:

\begin{itemize}
  \item Dacă este 0, $x$ este un element nou întîlnit.
  \item Dacă este 1, $x$ este 8.
  \item Dacă este 2, $x$ este 13.
  \item Dacă este 3, $x$ este 20.
\end{itemize}

Așadar, deși frecvențele sînt mari, putem ști exact ce fel de element este $x$ întrucît avem mai puțin de $k$ elemente distincte.

\subsubsection*{Cazul cu multe valori distincte}

Dacă numărul de elemente distincte este cel puțin $k$, vom folosi mai multe interogări, dar cum? Vom partiționa lista de frecvențe în $k$ grupe cît mai egale. Apoi vom construi o interogare astfel încît, în funcție de răspuns, să putem restrînge căutarea la una dintre grupe. De exemplu, fie $k=5$ și următoarea listă de frecvențe (inclusiv 0):

\begin{verbatim}
0 1 2 5 8 10 12 15 20 23 25 28 33
\end{verbatim}

Vom partiționa elementele în grupele [0 1 2], [5 8 10], [12 15 20], [23 25] și [28 33]. Vom construi interogarea astfel:

\begin{itemize}
  \item Elementele cu frecvențe 0, 1, 2 nu le luăm.
  \item Elementele cu frecvențe 5, 8, 10 le luăm o dată.
  \item Elementele cu frecvență 12, 15, 20 le luăm de 2 ori.
  \item Elementele cu frecvență 23, 25 le luăm de 3 ori.
  \item Elementele cu frecvență 28, 33 le luăm de 4 ori.
  \item Adăugăm și poziția curentă.
\end{itemize}

Acum, dacă răspunsul este 4, am redus posibilitățile la două: o frecvență de 28 a devenit 29 sau o frecvență de 33 a devenit 34. Facem a doua (și ultima) interogare pentru a diferenția aceste cazuri:

\begin{itemize}
  \item Elementele cu frecvență 28 nu le luăm.
  \item Elementele cu frecvență 33 le luăm o dată.
\end{itemize}

Acum, dacă răspunsul este 0, atunci explicația este că un element de frecvență 28 a crescut la 29. Altfel răspunsul este 1, ceea ce înseamnă că un element de frecvență 33 a crescut la 34.

\subsubsection*{Localizarea elementelor cu o frecvență cerută}

Mai este un ultim aspect, de care eu nu m-am prins pînă nu am citit \href{https://kilonova.ro/submissions/761967}{soluția aceasta}. Ca să adăugăm la interogare, de exemplu, elementele cu frecvențele 28 și 33 de 4 ori, trebuie să cunoaștem pozițiile aparițiilor acelor elemente. Dar cum? Din nou, nu este practic să încercăm să construim lista de poziții care conțin elemente egale.

Soluția este să notăm, pe fiecare poziție procesată, concluzia interogărilor: frecvența la care a ajuns acel element. Dacă un element a ajuns la frecvență 28, el a ajuns cîndva anterior și la frecvențele 1, 2, \dots, 27. Deci, pentru prima interogare de mai sus, putem proceda astfel:

\begin{itemize}
  \item Tipărim toate pozițiile pe care am ajuns la frecvență 5.
  \item Tipărim toate pozițiile pe care am ajuns la frecvență 12.
  \item Tipărim toate pozițiile pe care am ajuns la frecvență 23.
  \item Tipărim toate pozițiile pe care am ajuns la frecvență 28.
\end{itemize}

Atunci obținem ce dorim, adică:

\begin{itemize}
  \item Pozițiile elementelor cu frecvență 0, 1 sau 2 nu vor fi tipărite.

  \item Pozițiile elementelor cu frecvență 5, 8 sau 10 vor fi tipărite o singură dată, respectiv la momentul la care ating frecvența 5.

  \item \dots

  \item Pozițiile elementelor cu frecvență 28 sau 33 vor fi tipărite de 4 ori, respectiv la momentele la care ating frecvențele 5, 12, 23 și 28.
\end{itemize}

\subsubsection*{Detaliu de implementare}

Evaluatorul ne impune să nu folosim memorie globală, ceea ce are sens, deoarece cu memorie globală funcțiile \ccode{descuie_grajdul} și \ccode{dreseaza_cal} ar putea comunica direct între ele, șuntînd funcția \ccode{intreaba_cal} și făcînd problema banală. Din păcate, evaluatorul este prea zelos. El interzice complet declararea de \ccode{struct}-uri sau clase (mă întreb de ce!). De aceea, implementarea anexată este lizibilă, dar nu se va compila. Am rescris \href{https://kilonova.ro/submissions/772576}{sursa} pentru a obține 100p, dar o consider ilizibilă (cel puțin după standardele mele).

\end{document}
